\chapter{Implementácia}

\section{Vývojové prostredie}
Na vývoj aplikácie som využil Visual Studio Code, ktorý poskytuje širokú škálu rozšírení a nástrojov pre vývoj v jazyku Dart a frameworku Flutter. Pre správu závislostí a buildovanie projektu som použil nástroj Flutter CLI, ktorý je súčasťou Flutter SDK.

\section{Verzovací systém}
Na správu zdrojového kódu som využil verzovací systém Git.

\section{Získavanie dát}
Programová ponuka slovenského skautingu sa nachádza primárne na webových stŕankach. Obsahom stránok je detailný popis odboriek, výziev, voľných programových modulov a najvyšších programových ocenení. Na získanie dát ohľadom programu som zvolil webscraping\ref{webscraping} so snahou uložiť tieto dáta do štruktúrovanej podoby, ktorú môžem následne využiť v mojej aplikácii.

Analýzou webových stránok slovenského skautingu som dospel k nasledujúcej stratégií. V prvom kroku som získal postupne všetky odborky, ktoré boli na stránkach rozdelené podľa vekových kategórií. V druhom kroku som sa pustil do výziev ktoré boli náročnejšie v dôsledku laxnejšej štruktúry čomu bolo treba upraviť selektory na extrakciu textu a dátový model do ktorého som ich ukladal. V poslednom kroku som extrahoval z web stránky najvyššie programové ocenenia. Voľné programové moduly neobsahovali takmer žiadnu spoločnú štruktúru naprieč sebou, z tohto dôvodu boli vynechané počas získavania dát.

Na samotnú činnosť procesu webscraping som zvolil knižnicu todo pre programovací jazyk Python. Táto knižnica obsahovala potrebné nástroje na sťahovanie HTML obsahu webstránok a potrebné CSS selektory na extrakciu žiadaného textu.

Získané dáta som ukladal do súborov vo formáte JSON. Ktoré sú následne použité v rámci aplikácie.

\section{Flutter}
Flutter je open-source framework na tvorbu multiplatformných aplikácií. Aplikkácia je tvorená stromom widgetov ktoré tvoria jednotlivé časti používateľského rozhrania.

\subsection{Asynchrónne časti}
Vzhľadom na to že aplikácia musí komunikovat so serverom a databázou, je potrebné implementovat asynchrónne časti ktoré zabezpečia správnu komunikáciu a synchronizáciu dát. Flutter poskytuje nástroje na správu asynchrónnych operácií, ako sú Future a Stream, ktoré umožňujú efektívne načítavanie a aktualizáciu dát bez blokovania používateľského rozhrania. Future sa používa pre jednorazové asynchrónne operácie, zatiaľ čo Stream je vhodný pre kontinuálne prúdy dát, ako sú aktualizácie z databázy v reálnom čase alebo zmeny v stave autentizácie používateľa.

\subsection{Databázové spojenie}
Na potreby backendu som zvolil databázu Firestore v rámci Firebase, jedná sa o NoSQL JSON dokumentovú databázu. Pre integráciu Flutteru s Firestore je k dispozícií SDK od Googlu. V rámci SDK sú okrem základných operácií k dispozícií aj podpora offline zápisov, ktoré sa zapíšu do cloudu po obnovení internetového pripojenia klienta.

\subsection{Denormalizácia databázy}
Firestore ako NoSQL JSON dokumentová databáza nepodporuje klasický join z SQL a každý join by musel byť robený na strane klienta po získaní všetkých potrebných dát.

Obmedzenia v rámci Firestore databázy sú v celkovej veľkosti dokumentu a v počte prečítaných a zapísaných dokumentov, nie v množstve prenesených dát. Z tohto dôvodu som prispel k denormalizácií dát aby som znížil počet prečítaných dokumentov.

\subsubsection{Kolekcie}


\subsection{Zabezpečenie databáze}
Firestore umožňuje priamy prístup k dátam z klienta. Z tohto dôvodu existujú pravidlá ktoré umožnia zabezpečiť prístup k databázy a požiadavky ktoré nie sú povolené budú zamietnuté.

Pri nastavení pravidiel som postupoval podľa prípadov použiťia a povoloval iba operácie, ktoré boli potrebné.



\subsection{Nastavenie témy}
Na unifikáciu témy aplikácie flutter poskytuje nástroj ThemeData, ktorý umožňuje definovat farby, fonty a ďalšie vizuálne prvky aplikácie. V rámci implementácie som vytvoril vlastnú tému, ktorá reflektuje vizuálnu identitu slovenského skautingu, vrátane použití farieb a typografie, ktoré sú v súlade s oficiálnymi materiálmi skautingu.
% add theme snippet


\section{Využitie nástrojov umelej inteligencie}
V rámci implementácie autor využil nástroje umelej inteligencie na vygenerovanie boilerplate kódu a na opravu bugov počas vývoja. Pri generovaní boilerplate kódu si nástroje umelej inteligencie poradili s vysokou úspešnosťou avšak stále bolo treba doladiť vygenerovaný kód aby spĺňal požiadavky projektu. Pri oprave bugov si viedli nástroje umelej inteligencie podstatne horšie a často generovali kód ktorý nefungoval alebo bol neefektívny. Analýza kódu pomocou umelej inteligencie nebola dostatočná a často neidentifikovala problém.

\section{Nasadenie aplikácie}



\subsection{Zverejnenie aplikácie}